% Standalone problem document, generated from the adjacent README.md.
% Compile with XeLaTeX. Mathematical content is maintained in Markdown.
\documentclass[11pt,a4paper]{article}
\usepackage[fontsize=10.5pt]{scrextend}
\usepackage[margin=22mm,headheight=15pt,headsep=9mm,footskip=11mm]{geometry}
\usepackage{fontspec}
\setmainfont{texgyrepagella}[Extension=.otf,UprightFont=*-regular,BoldFont=*-bold,ItalicFont=*-italic,BoldItalicFont=*-bolditalic]
\setsansfont{texgyreheros}[Extension=.otf,UprightFont=*-regular,BoldFont=*-bold,ItalicFont=*-italic,BoldItalicFont=*-bolditalic]
\setmonofont{texgyrecursor}[Extension=.otf,UprightFont=*-regular,BoldFont=*-bold,ItalicFont=*-italic,BoldItalicFont=*-bolditalic,Scale=MatchLowercase]
\usepackage{amsmath,amssymb}
\usepackage{unicode-math}
\setmathfont{texgyrepagella-math.otf}
\usepackage{xcolor,microtype,enumitem,needspace,fancyhdr}
\usepackage{bookmark}
\definecolor{ink}{HTML}{17344B}
\definecolor{accent}{HTML}{197D80}
\definecolor{muted}{HTML}{596773}
\hypersetup{colorlinks=true,linkcolor=accent,urlcolor=accent,pdftitle={KE-01 - Exploit
spectral outliers without losing input
sparsity},pdfauthor={Open Problems in Numerical Linear Algebra}}
\urlstyle{same}
\setlength{\parindent}{0pt}
\setlength{\parskip}{5pt plus 1pt minus 1pt}
\linespread{1.04}
\setlength{\emergencystretch}{3em}
\widowpenalty=10000
\clubpenalty=10000
\displaywidowpenalty=10000
\allowdisplaybreaks[1]
\setlist{nosep,leftmargin=1.5em}
\providecommand{\tightlist}{\setlength{\itemsep}{0pt}\setlength{\parskip}{0pt}}
\providecommand{\pandocbounded}[1]{#1}
\pagestyle{fancy}
\fancyhf{}
\fancyhead[L]{\sffamily\footnotesize\color{muted}OPEN PROBLEMS IN NUMERICAL LINEAR ALGEBRA}
\fancyhead[R]{\sffamily\bfseries\footnotesize\color{accent}KE-01}
\fancyfoot[L]{\parbox[b]{0.9\textwidth}{\sffamily\fontsize{8}{10}\selectfont\color{muted}Editorial ratings; a bounded search cannot certify that a problem remains unsolved.\\Literature check: 2026-09-08}}
\fancyfoot[R]{\sffamily\footnotesize\color{muted}\thepage}
\renewcommand{\headrulewidth}{0.3pt}
\renewcommand{\headrule}{\hbox to\headwidth{\color{accent}\leaders\hrule height \headrulewidth\hfill}}
\makeatletter
\renewcommand{\section}{\Needspace{5\baselineskip}\@startsection{section}{1}{0pt}{12pt plus 2pt minus 2pt}{4pt}{\sffamily\bfseries\large\color{ink}}}
\renewcommand{\subsection}{\Needspace{4\baselineskip}\@startsection{subsection}{2}{0pt}{9pt plus 1pt minus 1pt}{3pt}{\sffamily\bfseries\normalsize\color{ink}}}
\makeatother
\setcounter{secnumdepth}{0}
\begin{document}
{\sffamily\small\color{muted}Linear systems and elimination\par}
\vspace{3pt}
{\raggedright\hyphenpenalty=10000\exhyphenpenalty=10000\sffamily\bfseries\fontsize{21}{24}\selectfont\color{ink}Exploit
spectral outliers without losing input sparsity\par}
\vspace{7pt}
{\sffamily\small\textbf{Difficulty:} challenging\quad\textbf{Importance:} interesting
to the community\par}
\vspace{4pt}
{\color{accent}\hrule height 0.8pt}
\vspace{6pt}
\textbf{Topic:} sparse linear systems; randomized preconditioning

\textbf{Status:} open; admitted after a bounded literature search found
no resolution.

\subsection{Context and notation}\label{context-and-notation}

All norms are Euclidean vector norms or their induced matrix norms.

\subsection{Problem statement}\label{problem-statement}

Let \(A\in\mathbb R^{n\times n}\) be nonsingular, with singular values
\(\sigma_1\geq\cdots\geq\sigma_n>0\), and write
\(N=\operatorname{nnz}(A)\) and \(\kappa_2(A)=\sigma_1/\sigma_n\). The
input consists of its nonzero entries and their indices, \(b\ne0\),
\(\varepsilon\in(0,1/2)\), \(k\in\{0,\ldots,n-1\}\), and
\(\kappa\geq1\), with the promise \(\sigma_{k+1}/\sigma_n\leq\kappa\).
Fix any exponent \(\omega_0>2\) for which square matrix multiplication
has an \(O(m^{\omega_0})\) arithmetic algorithm.

Does a randomized algorithm exist that, for every such input, produces
\(\widehat x\) satisfying

\[
\|A\widehat x-b\|_2\leq\varepsilon\|b\|_2
\]

with probability at least \(0.99\), using at most

\[
C\bigl(k^{\omega_0}+N\kappa\bigr)
\bigl[1+\log(n\kappa_2(A)/\varepsilon)\bigr]^q
\]

exact arithmetic operations? Here \(C,q\) may depend on the chosen
multiplication algorithm, but not on the input. Direct entry access is
allowed. The display makes the workshop's suppressed logarithms and
accuracy convention explicit; \(\omega_0\) avoids assuming that the
infimum defining the multiplication exponent is attained.

\subsection{References}\label{references}

Amsel et al.,
\href{https://arxiv.org/html/2602.05394v3\#S2.SS4}{\emph{Linear Systems
and Eigenvalue Problems}}, Definition 2.4 and Problem 2.5. Dereziński
and Sidford,
\href{https://doi.org/10.1137/1.9781611978971.37}{\emph{Approaching
Optimality for Solving Dense Linear Systems with Low-Rank Structure}},
SODA 2026, pp.~925--938; \href{https://arxiv.org/html/2507.11724}{full
manuscript}, Theorems 3 and 29, gives the dense-input predecessor.

\subsection{Status check}\label{status-check}

Searches for \textit{sparse outlying singular values solver 2026}
and \textit{spectral outliers linear systems 2026} found the newer
Liu, Nguyen, Peng, and Yang
\href{https://yangpliu.github.io/pdf/faster-solvers-sparse-spectral-outliers.pdf}{\emph{Faster
Solvers for Sparse Systems with Large Spectral Outliers}}. Its Theorem
1.1 assumes an explicitly given factor \(B\) in \(B^TBx=b\), bounded
spectral tail, polynomial conditioning, and short rational input. For an
\(m\times d\) factor with at most \(s\) nonzeros per row, its bit bound
includes \(m^{o(1)}\widetilde O(ms+\sqrt{d/k}\,k^2+k^{\omega+\eta})\).
These input restrictions and extra terms do not establish the requested
general bound. The workshop's August update records this as partial
progress on Problem 2.9.
\end{document}
